\section{Projects for Chapter 4: Conditional Probability, Independence, and Bayes' Theorem}

Conditional probability is one of the strongest candidates for visualisation because conditioning changes the reference population. The denominator should visibly change when new information is imposed.

\subsection*{Project: conditioning changes the space}

\begin{examplebox}
\textbf{Prompt to give an AI coding assistant}

\begin{quote}\small
Create an interactive finite probability space displayed as a grid of weighted elementary outcomes. Let the user define events \(A\) and \(B\). Show \(P(A)\), \(P(B)\), \(P(A\cap B)\), \(P(A\mid B)\), and \(P(B\mid A)\).

When the user activates conditioning on \(B\), visually fade all outcomes outside \(B\), renormalise the remaining weights, and show that
\[
P(A\mid B)=\frac{P(A\cap B)}{P(B)}.
\]
Include presets for disjoint events, nested events, independent events, and dependent events. Display whether the equality \(P(A\cap B)=P(A)P(B)\) holds.
\end{quote}
\end{examplebox}

\subsection*{Project: Bayes diagnostic-test laboratory}

This project should use prevalence, sensitivity, and specificity as adjustable parameters. A tree diagram and an icon array should display the same model.

\begin{examplebox}
\textbf{Prompt to give an AI coding assistant}

\begin{quote}\small
Build a Bayes theorem laboratory for a diagnostic test. Provide sliders for prevalence, sensitivity, specificity, and population size. Show a probability tree, an expected-count table, and an icon array. Calculate the positive predictive value and negative predictive value from exact probabilities.

Highlight the difference between \(P(\text{positive}\mid\text{condition})\) and \(P(\text{condition}\mid\text{positive})\). Add a simulation mode that generates individuals and compares empirical proportions with the exact Bayes calculation. Include a button that changes only prevalence so the student can see why predictive values change even when the test characteristics remain fixed.
\end{quote}
\end{examplebox}

An extension may display a partition \(B_1,\ldots,B_k\) and animate the law of total probability before reversing the reasoning with Bayes' theorem.

\begin{summarybox}
The key visual principle is that conditioning restricts the space and then rescales probability inside the restricted space.
\end{summarybox}
